% =====================================================================
\section{Exercise 3 -- Conditional quantile curves of the NM copula}\label{sec:ex3}
% =====================================================================

% ---------------------------------------------------------------
\subsection{Specification}
% ---------------------------------------------------------------

The Normal-Mixture (NM) copula is the convex combination of two Gaussian
copulas with different correlations:
\[
C^{NM}(u_1,u_2)=\pi\,C^G_{\rho_1}(u_1,u_2)+(1-\pi)\,C^G_{\rho_2}(u_1,u_2),
\quad 0<\pi<1.
\]
Since the conditional CDF is linear in the components, the conditional NM CDF
inherits the same structure:
\begin{equation}\label{eq:nm-cond}
C^{NM}(u_1\given u_2)
=\pi\,\Phi\!\Bigl(\tfrac{\Phi^{-1}(u_1)-\rho_1\Phi^{-1}(u_2)}{\sqrt{1-\rho_1^2}}\Bigr)
+(1-\pi)\,\Phi\!\Bigl(\tfrac{\Phi^{-1}(u_1)-\rho_2\Phi^{-1}(u_2)}{\sqrt{1-\rho_2^2}}\Bigr).
\end{equation}
Parameters used: $\pi=\ExThreePi,\;\rho_1=\ExThreeRhoA,\;\rho_2=\ExThreeRhoB$.

% ---------------------------------------------------------------
\subsection{Definition of $q$-quantile curves}
% ---------------------------------------------------------------

For each conditioning value $u_2\in(0,1)$ and each level $q\in(0,1)$ the
$q$-quantile curve is the locus
\[
\bigl\{(u_2,u_1)\in(0,1)^2 : C^{NM}(u_1\given u_2)=q\bigr\}.
\]
Because $u_1\mapsto C^{NM}(u_1\given u_2)$ is a CDF on $(0,1)$ (hence
\emph{strictly increasing} from 0 to 1), the equation has a unique root for
every $q\in(0,1)$. We solve it by Brent's method on $[\varepsilon,1-\varepsilon]$
with $\varepsilon=10^{-4}$ in
\code{exercise3.nm\_copula.nm\_quantile\_u1}.

\begin{remark}[Why we invert in $u_1$ and not in $u_2$]
The reading of the problem statement asks to obtain $u_2$
given $u_1$. However, with $\pi=\ExThreePi$, $\rho_1=\ExThreeRhoA$, $\rho_2=\ExThreeRhoB$ the map
$u_2\mapsto C^{NM}(u_1\given u_2)$ has range bounded inside $(0.3,0.7)$ for a
fixed $u_1$ in the interior, so no $u_2$-solution exists for $q\in\{0.05,0.95\}$.
Consistently with the slides we invert in $u_1$, which has range $(0,1)$ and a unique
solution for every $q$.
\end{remark}

% ---------------------------------------------------------------
\subsection{Numerical results}
% ---------------------------------------------------------------

Table~\ref{tab:3-u1} reports $\hat u_1(u_2;q)$ on a representative grid;
all roots verify $|C^{NM}(\hat u_1\given u_2)-q|<10^{-10}$.
Figure~\ref{fig:3} shows the curves on the $(u_2,u_1)$ unit square (left)
and on the $\mathcal{N}(0,1)$ return scale $(r_2,r_1)=(\Phi^{-1}u_2,\Phi^{-1}u_1)$
(right).

\begin{table}[H]\centering\small
\caption{$\hat u_1$ such that $C^{NM}(\hat u_1\given u_2)=q$.}\label{tab:3-u1}
\begin{tabular}{cccccc}
\toprule
\ExThreeQuantileHeader
\midrule
\ExThreeQuantileBody
\bottomrule
\end{tabular}
\end{table}

\begin{figure}[H]\centering
\includegraphics[width=\linewidth]{fig_3_nm_quantile_curves.png}
\caption{NM-copula $q$-quantile curves on the $(u_2,u_1)$ unit square (left)
and on the $\mathcal{N}(0,1)$ return scale (right). Parameters
$\pi=\ExThreePi$, $\rho_1=\ExThreeRhoA$, $\rho_2=\ExThreeRhoB$.}
\label{fig:3}
\end{figure}

\paragraph{Intuition.}
The median curve ($q=0.5$) is symmetric about $u_2=0.5$ because, by
construction, the conditional median of a continuous copula passes through
$(0.5,0.5)$. The outer quantiles are visibly asymmetric: in the lower region ($u_2$ small) the
$q=0.05$ branch is close to zero by the negative-correlation ($\rho_1=\ExThreeRhoA$, weight $\ExThreePi$);
in the positive region ($u_2$ large) the $q=0.05$ quantile is lifted because the positive-correlation
($\rho_2=\ExThreeRhoB$) component dominates. The mixture therefore generates a non-monotone,
``bow-shaped'' tail behaviour.
